\documentclass[10pt,a4paper]{amsart}
\usepackage{amsmath,amssymb,amsthm,graphicx}
% Retained minimal preamble for the Stacks Project chapter more-morphisms.tex (More on Morphisms).
% Verbatim excerpt of the chapter preamble: Xy-pic setup (native preamble lines 14-18) and the native
% theorem declarations with the mathematical macros (native preamble lines 154-217). Copied mechanically;
% nothing added, removed or redefined. The wrapper class and the amsmath/amssymb/amsthm/graphicx base
% packages are supplied by the assembler, so the class/externaldocument part of the native preamble is omitted.
% For commutative diagrams we use Xy-pic
\usepackage[all]{xy}

% We use 2cell for 2-commutative diagrams.
\xyoption{2cell}
% Theorem environments.
%
\theoremstyle{plain}
\newtheorem{theorem}[subsection]{Theorem}
\newtheorem{proposition}[subsection]{Proposition}
\newtheorem{lemma}[subsection]{Lemma}

\theoremstyle{definition}
\newtheorem{definition}[subsection]{Definition}
\newtheorem{example}[subsection]{Example}
\newtheorem{exercise}[subsection]{Exercise}
\newtheorem{situation}[subsection]{Situation}

\theoremstyle{remark}
\newtheorem{remark}[subsection]{Remark}
\newtheorem{remarks}[subsection]{Remarks}

\numberwithin{equation}{subsection}

% Macros
%
\def\lim{\mathop{\mathrm{lim}}\nolimits}
\def\colim{\mathop{\mathrm{colim}}\nolimits}
\def\Spec{\mathop{\mathrm{Spec}}}
\def\Hom{\mathop{\mathrm{Hom}}\nolimits}
\def\Ext{\mathop{\mathrm{Ext}}\nolimits}
\def\SheafHom{\mathop{\mathcal{H}\!\mathit{om}}\nolimits}
\def\SheafExt{\mathop{\mathcal{E}\!\mathit{xt}}\nolimits}
\def\Sch{\mathit{Sch}}
\def\Mor{\mathop{\mathrm{Mor}}\nolimits}
\def\Ob{\mathop{\mathrm{Ob}}\nolimits}
\def\Sh{\mathop{\mathit{Sh}}\nolimits}
\def\NL{\mathop{N\!L}\nolimits}
\def\CH{\mathop{\mathrm{CH}}\nolimits}
\def\proetale{{pro\text{-}\acute{e}tale}}
\def\etale{{\acute{e}tale}}
\def\QCoh{\mathit{QCoh}}
\def\Ker{\mathop{\mathrm{Ker}}}
\def\Im{\mathop{\mathrm{Im}}}
\def\Coker{\mathop{\mathrm{Coker}}}
\def\Coim{\mathop{\mathrm{Coim}}}

% Boxtimes
%
\DeclareMathSymbol{\boxtimes}{\mathbin}{AMSa}{"02}

%
% Macros for moduli stacks/spaces
%
\def\QCohstack{\mathcal{QC}\!\mathit{oh}}
\def\Cohstack{\mathcal{C}\!\mathit{oh}}
\def\Spacesstack{\mathcal{S}\!\mathit{paces}}
\def\Quotfunctor{\mathrm{Quot}}
\def\Hilbfunctor{\mathrm{Hilb}}
\def\Curvesstack{\mathcal{C}\!\mathit{urves}}
\def\Polarizedstack{\mathcal{P}\!\mathit{olarized}}
\def\Complexesstack{\mathcal{C}\!\mathit{omplexes}}
% \Pic is the operator that assigns to X its picard group, usage \Pic(X)
% \Picardstack_{X/B} denotes the Picard stack of X over B
% \Picardfunctor_{X/B} denotes the Picard functor of X over B
\def\Pic{\mathop{\mathrm{Pic}}\nolimits}
\def\Picardstack{\mathcal{P}\!\mathit{ic}}
\def\Picardfunctor{\mathrm{Pic}}
\def\Deformationcategory{\mathcal{D}\!\mathit{ef}}
\usepackage{cleveref}
\crefname{originaldefinition}{original definition}{original definitions}
\Crefname{originaldefinition}{Original definition}{Original definitions}
\makeatletter
\newlabel{definition-weighting}{{0F3A}{source}{}{sourcepointer-1}{}}
\newlabel{definition-weighting@cref}{{[originaldefinition][0][]0F3A}{[1][0][]source}{}{}{}}
\makeatother
\crefname{originallemma}{original lemma}{original lemmas}
\Crefname{originallemma}{Original lemma}{Original lemmas}
\makeatletter
\newlabel{lemma-etale-makes-quasi-finite-finite-at-point}{{02LK}{source}{}{sourcepointer-2}{}}
\newlabel{lemma-etale-makes-quasi-finite-finite-at-point@cref}{{[originallemma][0][]02LK}{[1][0][]source}{}{}{}}
\makeatother
\title{A weighting vanishing over one fibre vanishes over every specialisation of that fibre}
\author{The Stacks Project Authors}
\date{stacks-project:more-morphisms.tex:local-snapshot-20261011; GNU FDL 1.2 or later}
\begin{document}
\maketitle
\noindent\textit{Editorial note.} Counted span is the author's own proof environment, source lines 22477-22499, immediately following the statement at 22468-22475.\par\medskip
\noindent\textit{Editorial note.} The proof cites one chapter lemma and the definition of a weighting by \textbackslash{}ref; both are navigation pointers into the same chapter and are not re-counted.\par\medskip
\noindent\textit{Editorial note.} No step of the proof is omitted by the author; the closing sentence is a genuine deduction from the displayed value of the integral at v'.\par\medskip
\noindent\textit{Allowed standard prerequisites (editorial).}\begin{itemize}
\item Definition of a weighting and of its integral. chapter definition definition-weighting, cited at proof line 22492. Supplies the local constancy of the integral of the weighting along a finite morphism over the base.
\item Quasi-finiteness becomes finiteness after an elementary etale neighbourhood. chapter lemma lemma-etale-makes-quasi-finite-finite-at-point, cited at proof line 22479. Produces the etale neighbourhood V -> Y, the open U of X\_V and the finite U -> V with a unique point over the marked point.
\item Openness of etale morphisms. standard Stacks fact used at proof line 22490. Produces a point v of V specialising to v' and mapping to the smaller point y of Y.
\item Separable degree of a residue field extension. the factor [kappa(u'):kappa(v')]\_s appearing at proof line 22497. Identifies the value of the integral at v' with the weighting of x' times a separable degree, which is nonzero so that vanishing of the integral forces w(x') = 0.
\end{itemize}\medskip
\ifdefined\hypertarget\hypertarget{sourcepointer-1}{}\fi
\noindent\textit{Source reference.} Original definition 0F3A: sources/stacks-more-morphisms/pinned-original-source.tex, source lines 22152--22152. Navigation pointer only: the counted proof cites this statement by \textbackslash{}ref as a standing input; the deduction the counted proof has to make is carried out inside the counted span, and no part of that deduction is delegated to this pointer.\par\medskip
\ifdefined\hypertarget\hypertarget{sourcepointer-2}{}\fi
\noindent\textit{Source reference.} Original lemma 02LK: sources/stacks-more-morphisms/pinned-original-source.tex, source lines 11924--11924. Navigation pointer only: the counted proof cites this statement by \textbackslash{}ref as a standing input; the deduction the counted proof has to make is carried out inside the counted span, and no part of that deduction is delegated to this pointer.\par\medskip
\begin{lemma}
\label{lemma-weighting-specialization}
Let $f : X \to Y$ be a locally quasi-finite morphism of schemes.
Let $w : X \to \mathbf{Z}$ be a weighting of $f$.
Let $y \leadsto y'$ be a specialization of points of $Y$.
Assume $w(x) = 0$ for all $x \in X$ with $f(x) = y$. Then
$w(x') = 0$ for all $x' \in X$ with $f(x') = y'$.
\end{lemma}

\begin{proof}
Let $x' \in X$ be a point mapping to $y'$. By
Lemma \ref{lemma-etale-makes-quasi-finite-finite-at-point}
we can choose a diagram
$$
\xymatrix{
U \ar[r] \ar[dr]_\pi & X_V \ar[d] \ar[r] & X \ar[d]^f \\
& V \ar[r]^\varphi & Y
}
$$
with $V \to Y$ \'etale, with a point $v' \in V$ mapping to $y'$,
with $U \subset Y_V$ open, such that $U \to V$ is finite, and
such that there exists an unique point $u' \in U$ mapping to $v'$
which moreover maps to $x'$ in $X$. Note that since $V \to Y$ is open,
there exists a point $v \in V$ specializing to $v'$ and mapping to $y$.
By Definition \ref{definition-weighting} the function
$\int_\pi (w \circ h)$ is locally constant where $h : U \to X$ is
the composition of the top horizontal arrows.
Since $\int_\pi (w \circ h)$ is zero in $v$ by
our assumption on $w$, we conclude that it is zero
at $v'$. However, the value at $v'$ is $w(x')[\kappa(u'):\kappa(v')]_s$
and we conclude.
\end{proof}

\end{document}
